\babbaru{BAB IV HASIL DAN PEMBAHASAN}{%
\petunjuk{Inti laporan: temuan dan interpretasinya (Pedoman TA 7.4). Revisi yang mengubah substansi bab ini tergolong revisi sedang atau mayor.}{\statusdua{\TIDAKPERLU}{\WAJIB}}
}

\subbagian{4.1 Hasil Penelitian}
\phantomsection\addcontentsline{toc}{section}{4.1 Hasil Penelitian}
\petunjuk{Sajikan hasil sesuai urutan tahapan metodologi dalam tabel, grafik, atau diagram yang bernomor, berjudul, dan dirujuk di narasi.}{}

\placeholder{[Tulis uraian dalam bentuk paragraf.]}\par

\subbagian{4.2 Analisis dan Pembahasan}
\phantomsection\addcontentsline{toc}{section}{4.2 Analisis dan Pembahasan}
\petunjuk{Jelaskan makna dan pola data, bandingkan dengan penelitian terdahulu (Bab 2.2), identifikasi faktor yang memengaruhi hasil, dan jawab setiap rumusan masalah. Jangan sekadar mengulang isi tabel.}{}

\placeholder{[Tulis uraian dalam bentuk paragraf.]}\par

\subbagian{4.3 Keterbatasan Penelitian}
\phantomsection\addcontentsline{toc}{section}{4.3 Keterbatasan Penelitian}
\petunjuk{Akui secara jujur keterbatasan metode, dataset, lingkungan eksperimen, dan cakupan analisis.}{}

\placeholder{[Tulis uraian dalam bentuk paragraf.]}\par
\clearpage
